\chapter{Inward Pointing Vectors in Manifolds with Corners}

Defining the notion of inward pointing vectors in manifolds with corners is substantially harder than in manifolds with boundary without corners. Let us first examine why this is. 

Here is a common definition for manifolds with boundary without corners:

\begin{defi}[\cite{Lee2006} p.\ 136]\label{defi:InwLee1}
    Let $M$ be a smooth manifold with boundary without corners. Let $p \in M$. Then $v \in \tanspace{p}{M} \setminus \tanspace{p}{\boundary M}$ is inward pointing if and only if there is an $\varepsilon > 0$ and a smooth curve $\gamma \colon [0, \varepsilon) \to M$ such that $\gamma (0) = p$ and $\gamma' (0) = v$.
\end{defi}

This definition cannot be used for manifolds with corners as in this case $\boundary M$ is not always a manifold and thus $\tanspace{p}{\boundary M}$ is not defined. 

The existence of the $\gamma \colon : [0, \varepsilon) \to M$ by itself characterises whether a vector is inward pointing or tangent to the boundary, a notion which we will call \emph{realisable}. 

So for manifolds with corners we either need a completely new definition that intrinsically excludes the vectors tangent to the boundary or we need a definition of realisability and a way to exclude the tangent vectors. 

We first tried the first option and came up with these candidates for definitions which we either came up with ourselves or found in the literature: 

\begin{defi}[Our idea inspired by Definition \ref{defi:InwLee1}]
    Let $M$ be a smooth manifold with corners. Let $p \in M$. Then $v \in \tanspace{p}{M}$ is inward pointing if and only if there is an $\varepsilon > 0$ and a smooth curve $\gamma \colon [0, \varepsilon) \to M$ such that $\gamma (0) = p$, $\gamma' (0) = v$ and $\gamma ((0, \varepsilon)) \subseteq \interior{M}$.
\end{defi}

Unfortunately, this definition doesn't exclude vectors tangent to the boundary: 

\begin{example}
    Take $M$ to be the upper half plane $\mathcal{H}$ in $\bR^2$ let $p$ be $(0, 0)$ and let $\gamma \colon [0, 1) \to \mathcal{H}, x \mapsto x^2$. Then $\gamma'(0) = (1, 0)$ and thus by this definition $(1, 0)$ would be inward pointing even though it is clearly tangent to the boundary.
\end{example}

\begin{figure}
\centering
\begin{tikzpicture}[scale=3]
    \draw (0, 0) parabola (1, 1);
    \draw (-1, 0) -- (1, 0);
    \node[label={[yshift=0.3](0, 0)}, circle, fill, inner sep=1.5pt] at (0, 0) {};
    \node at (0.6, 0.5) {$\gamma$};
    %\fill (-1, 1) -- (-1, 0) -- (1, 1) -- (0, 1)
\end{tikzpicture}
\end{figure}

\begin{defi}
    Let $M$ be a smooth manifold with corners. Let $p \in M$. Then $v \in \tanspace{p}{M}$ is inward pointing if and only if $\direcderiv{v}f (p) \ge 0$ for every non-negative $f \in \smoothfunc{M}$ that is zero on $\boundary{M}$. 
\end{defi}